\section{Final Calculation Results}

\begin{table}[h]
\centering
\caption{Comparison of Programming Languages by Metrics}
\label{tab:languages_metrics}
\begin{tabularx}{\textwidth}{|l|X|X|X|X|X|X|X|X|}
\hline
\textbf{Language} & 
\textbf{TSSI} & 
\textbf{AE} & 
\textbf{ISRI} & 
\textbf{CRFI} & 
\textbf{ORCI} & 
\textbf{PDEI} & 
\textbf{SCSI} & 
\textbf{FSVI} \\
\hline

C\# & 0.7833 & 0.3875 & 0.9083 & 0.5 & 0.7433 & 0.685 & 0.25 & 0.275 \\ \hline
Java & 0.5 & 0.3643 & 0.731 & 0.425 & 0.6917 & 0.545 & 0.25 & 0.275 \\ \hline
Smalltalk & 0 & 0.0286 & 0.45 & 0.225 & 0.5417 & 0.07 & 0 & 0 \\ \hline
C++ & 0.3 & 0.3774 & 0.3907 & 0.425 & 0.3417 & 0.37 & 0.1667 & 0.275 \\ \hline
C & 0 & 0.0571 & 0 & 0.075 & 0.15 & 0.1 & 0.1667 & 0 \\ \hline
Golang & 0.6334 & 0.4321 & 0.7167 & 0.25 & 0.5917 & 0.245 & 0.3 & 0.275 \\ \hline
Scala & 0.8 & 0.3381 & 0.8167 & 0.6875 & 0.6917 & 0.62 & 0.4 & 0.275 \\ \hline
JavaScript & 0.2 & 0.3071 & 0.4917 & 0.075 & 0.925 & 0.07 & 0.1 & 0.275 \\ \hline
Python & 0.8583 & 0.5207 & 0.3417 & 0.5125 & 0.85 & 0.07 & 0.1 & 0.275 \\ \hline
Ruby & 0.15 & 0.2107 & 0.4917 & 0.5125 & 0.925 & 0.07 & 0 & 0.275 \\ \hline
Zonnon & 0.55 & 0.8232 & 0.7667 & 0.25 & 0.4417 & 0.445 & 0.3667 & 0.575 \\ \hline

\end{tabularx}
\end{table}

\subsection{Comparison summary: C\#}

The obtained results demonstrate that C\# provides one of the most structurally disciplined object-oriented models among mainstream industrial languages, particularly in the areas of inheritance management and type-system safety. The high values of $ISRI(\text{C\#}) \approx 0.9083$ and $TSSI(\text{C\#}) \approx 0.7833$ indicate that the language strongly separates inheritance management concerns from subtype relations while simultaneously maintaining a formally constrained nominal type system with explicit variance annotations. This result is consistent with the design goals of the .NET type system, where generic variance and interface-based abstraction were introduced specifically to improve type safety and behavioral substitutability \cite{pierce2002types}. In particular, the use of explicit variance modifiers (\texttt{in}/\texttt{out}) and deterministic inheritance semantics reduces ambiguity in polymorphic hierarchies and mitigates the fragile base class problem through mandatory override rules and sealed constructs. Such characteristics correspond to the observations of Cardelli and Cook that inheritance and subtyping should remain conceptually distinct relations in robust object-oriented systems \cite{cardelli1985understanding,cook1989inheritance}. The relatively high $ORCI(\text{C\#}) \approx 0.7433$ additionally reflects the flexibility of the runtime object model and the availability of controlled reflective mechanisms through the Common Language Runtime.

At the same time, the lower scores in $AE(\text{C\#}) = 0.3875$, $SCSI(\text{C\#}) = 0.25$, and $FSVI(\text{C\#}) = 0.275$ reveal several structural limitations of the language from the perspective of strict object-oriented correctness and formal verifiability. Despite providing multiple visibility modifiers and interface abstractions, C\# does not enforce strong module isolation and allows extensive runtime reflection, which weakens encapsulation guarantees. This observation corresponds to the critique of unrestricted reflective systems discussed by Bracha and Ungar, who note that reflection often compromises abstraction barriers when not separated through dedicated metaobject protocols \cite{bracha2004mirrors}. Furthermore, mutability-by-default semantics and the absence of ownership-based aliasing restrictions significantly reduce the language's static guarantees regarding concurrency safety, despite the availability of higher-level asynchronous constructs. Research on ownership types and alias control emphasizes that unrestricted mutable sharing substantially complicates reasoning about correctness in concurrent systems \cite{clarke1998ownership}. Finally, the low $FSVI$ value reflects the absence of native Design by Contract mechanisms in the language specification, since preconditions, postconditions, and invariants rely primarily on external libraries or runtime assertions rather than first-class language constructs. As argued by Meyer, the absence of specification-level contracts limits the possibility of integrating correctness reasoning directly into the object model itself \cite{meyer2002applying}. Overall, the results indicate that C\# prioritizes practical industrial flexibility and interoperability over strict encapsulation, formal verification, and state safety guarantees.

\subsection{Comparison summary: Java}

The obtained metric values indicate that Java provides a relatively mature and formally constrained object model, but one that remains strongly oriented toward classical class-based object-oriented programming. The value of $TSSI(\text{Java}) = 0.5$ reflects the hybrid nature of Java’s type system: although the language separates implementation inheritance from subtype polymorphism through interfaces, its subtype model remains fundamentally nominal. This corresponds to the observations of Cardelli, who distinguishes nominal inheritance systems from structurally defined subtyping systems and notes that nominal models constrain substitutability to explicitly declared hierarchies \cite{cardelli1985understanding}. Java’s generic variance model is also intentionally conservative due to type-erasure compatibility and wildcard restrictions, which partially limits formal expressiveness despite preserving type safety \cite{pierce2002types}. At the same time, the relatively high value of $ISRI(\text{Java}) \approx 0.731$ demonstrates that the language specification places strong emphasis on deterministic inheritance semantics and hierarchy control. Java avoids multiple implementation inheritance, defines deterministic method resolution, and includes mechanisms such as \texttt{final}, explicit \texttt{override}-style semantics, and sealed hierarchies, which collectively reduce ambiguity and mitigate fragile base class effects. Such restrictions are consistent with the broader tendency in object-oriented language design to prioritize predictability and maintainability over unrestricted inheritance flexibility \cite{cook1989inheritance, scharli2003traits}. 

The remaining metrics reveal several limitations of Java from the perspective of modern object-oriented expressiveness. The relatively low values of $AE(\text{Java}) = 0.3643$ and $CRFI(\text{Java}) = 0.425$ indicate that Java provides only moderate support for abstraction-oriented composition mechanisms beyond conventional interfaces. Although interfaces with default methods partially increase behavioral composability, the language lacks native trait systems, structured delegation, and explicit conflict-resolution mechanisms characteristic of more composition-oriented models \cite{scharli2003traits}. The value of $ORCI(\text{Java}) \approx 0.6917$ confirms that Java preserves a strong and consistent reference-based object identity model, while also supporting several object instantiation paradigms. However, metaprogramming capabilities remain intentionally restricted due to the controlled reflection model of the JVM. This design corresponds to the argument of Bracha and Ungar that unrestricted reflection often conflicts with encapsulation guarantees and system safety \cite{bracha2004mirrors}. The medium value of $PDEI(\text{Java}) = 0.545$ further demonstrates that Java supports only classical single dispatch, without native multiple dispatch or predicate dispatch mechanisms, thereby limiting extensibility of polymorphic behavior compared to more expressive object systems \cite{castagna1995covariance}. Finally, the lowest values were observed for $SCSI(\text{Java}) = 0.25$ and $FSVI(\text{Java}) = 0.275$. These results show that Java relies predominantly on programmer discipline for mutability and concurrency safety and does not provide native Design-by-Contract constructs. While runtime assertion libraries and static analyzers partially compensate for these limitations, formal specification mechanisms are not embedded into the core language semantics, unlike approaches proposed in Eiffel and contract-oriented systems \cite{meyer2002applying}. Overall, the results characterize Java as a conservative and industrially stable object-oriented language whose design prioritizes compatibility, predictability, and controlled extensibility over maximal flexibility of object composition and formal behavioral specification.

\subsection{Comparison summary: Smalltalk}

The obtained metric values demonstrate that Smalltalk embodies one of the earliest and most conceptually pure object-oriented models, but at the same time lacks many of the formal safety and modularity mechanisms that later became central in statically analyzed object-oriented languages. The zero value of the TSSI metric indicates the absence of a formally specified type and subtyping system. This result is consistent with the classical description of Smalltalk as a dynamically typed pure object-oriented language in which polymorphism is based on message passing rather than static subtype relations \cite{goldberg1983smalltalk}. Since the language specification does not distinguish nominal and structural subtyping and provides no formal variance model, type compatibility is resolved entirely at runtime. Such an approach significantly increases flexibility but removes compile-time guarantees discussed in type-theoretic studies by Cardelli and Pierce \cite{cardelli1985understanding, pierce2002types}. Similarly, the extremely low AE and PDEI values reflect the historically minimalistic encapsulation model of Smalltalk. The language traditionally exposes most object structure through reflection and image-level introspection, while method dispatch relies exclusively on dynamic single dispatch without static exhaustiveness guarantees or devirtualization mechanisms.

At the same time, several metrics reveal important architectural strengths of the language. The ORCI value is comparatively high due to the consistency of reference semantics and the highly reflective metaobject environment. Smalltalk was historically one of the first languages to treat the runtime image, classes, and metaclasses as uniformly accessible objects, which corresponds to the reflective object model described in \cite{kiczales1991art}. The ISRI score also demonstrates that despite the lack of formal restrictions, the inheritance model itself remains relatively simple and deterministic because Smalltalk uses only single inheritance with a well-defined message lookup order. However, the low CRFI value shows that the language predates modern composition-oriented reuse mechanisms such as traits, mixins, and interface default implementations. This limitation became one of the motivations for later research on traits in dynamically typed object-oriented systems \cite{scharli2003traits}. Finally, the zero values for SCSI and FSVI indicate that the language specification contains neither static concurrency guarantees nor native support for contracts and formal verification. Consequently, the obtained results characterize Smalltalk as a historically foundational but weakly formalized object-oriented system whose design prioritizes dynamism, reflection, and conceptual uniformity over static correctness, modular isolation, and verifiability.

\subsection{Comparison summary: C++}

The obtained metric values demonstrate that the object model of \texttt{C++} remains strongly oriented toward classical implementation inheritance and low-level memory control rather than toward formally safe and semantically isolated object-oriented abstractions. The low value of the Type System and Subtyping Soundness Index ($TSSI = 0.3$) indicates that subtyping in \texttt{C++} is almost entirely derived from inheritance relations, while the language lacks formal separation between subtype compatibility and implementation reuse. This observation is consistent with the analysis of Cardelli, who emphasized that inheritance and subtyping represent conceptually different mechanisms and should not be treated as equivalent relations \cite{cardelli1985understanding}. Although \texttt{C++} provides nominal typing and generic mechanisms through templates, variance rules are not formally integrated into the type system specification, which reduces predictability of generic substitutability. Furthermore, the relatively low values of the Abstraction and Encapsulation score ($AE \approx 0.3774$) and the Inheritance Semantics Robustness Index ($ISRI \approx 0.3907$) reflect the fact that the language prioritizes flexibility over encapsulation guarantees. In particular, unrestricted reflection is absent, yet direct memory manipulation and unsafe casts allow bypassing abstraction boundaries, which weakens encapsulation integrity. Multiple inheritance additionally increases hierarchy complexity and introduces ambiguity in method resolution, a problem extensively discussed in studies on inheritance semantics and trait-based composition \cite{cook1989inheritance,scharli2003traits}.

The metrics related to composition, state safety, and formal verification reveal even more substantial limitations of the \texttt{C++} object model from the perspective of robust object-oriented design. The Composition and Reuse Flexibility Index ($CRFI = 0.425$) demonstrates that although interfaces can be composed flexibly through multiple inheritance and abstract classes, the language lacks native trait systems and structured delegation mechanisms. As noted by Schärli et al., traits were specifically introduced to avoid the semantic fragility and ambiguity caused by inheritance-based reuse \cite{scharli2003traits}. The low values of the Object Representation and Creation Index ($ORCI \approx 0.3417$) and the Polymorphism and Dispatch Efficiency Index ($PDEI = 0.37$) further indicate that \texttt{C++} retains a predominantly class-centric and manually controlled runtime model with limited metaobject extensibility and only single-dispatch polymorphism. Most critically, the State and Concurrency Safety Index reaches the lowest result among all evaluated dimensions ($SCSI \approx 0.1667$), demonstrating the absence of immutable-by-default semantics, ownership enforcement, and compile-time guarantees of data-race freedom. This result directly corresponds to the observations of Clarke et al. regarding the importance of ownership and aliasing restrictions for safe concurrent object manipulation \cite{clarke1998ownership}. Finally, the Formal Specification and Verifiability Index ($FSVI = 0.275$) confirms that Design by Contract principles are not represented as first-class language constructs in \texttt{C++}, correctness constraints rely mainly on external libraries and runtime assertions rather than on formally verifiable specifications, contrary to the approach proposed by Meyer \cite{meyer2002applying}. Overall, the resulting metric profile characterizes \texttt{C++} as a language with extensive expressive power and implementation flexibility, but with comparatively weak formal guarantees for safe, verifiable, and semantically isolated object-oriented programming.

\subsection{Comparison summary: C}

The obtained metric values demonstrate that the C programming language provides almost no native support for object-oriented abstractions and therefore serves primarily as a procedural systems language rather than an object-oriented environment. The value $TSSI(C)=0$ reflects the absence of a formal subtype system, inheritance semantics, variance rules, and separation between inheritance and substitutability. This result is consistent with the classical distinction between procedural and object-oriented type systems described by Cardelli and Wegner, who emphasize that subtype polymorphism and inheritance relations are fundamental components of object-oriented languages \cite{cardelli1985understanding}. Similarly, the values $ISRI(C)=0$ and $FSVI(C)=0$ indicate that the language specification contains no formal inheritance model, no hierarchy management mechanisms, and no native support for contracts or formal behavioral verification. Such characteristics correspond to the historical design goals of C as a portable low-level systems language prioritizing direct memory access and minimal runtime overhead rather than semantic safety or formal abstraction mechanisms.

The low values of $AE(C)\approx0.0571$, $CRFI(C)=0.075$, and $PDEI(C)=0.1$ further demonstrate that encapsulation, compositional reuse, and polymorphic dispatch are not represented as first-class language concepts in C. Encapsulation is limited to translation-unit visibility through the \texttt{static} keyword, while module isolation and interface abstraction are absent at the language level. As noted by Parnas, effective modular decomposition requires explicit information hiding boundaries \cite{parnas1972criteria}, which are not formally enforced in C. The value $ORCI(C)=0.15$ indicates that although the language supports several allocation paradigms (automatic and dynamic allocation), it lacks a unified object model and reflective meta-level capabilities. The result $SCSI(C)\approx0.1667$ also highlights that mutability is unrestricted by default and concurrency correctness depends entirely on programmer discipline, which corresponds to observations in ownership-type research emphasizing that unrestricted aliasing significantly complicates safe concurrent programming \cite{clarke1998ownership}. Overall, the obtained metrics confirm that C provides only primitive mechanisms from which object-oriented abstractions may be manually emulated, but the language itself does not formally support the semantic guarantees expected from modern object-oriented systems.

\subsection{Comparison summary: Golang}

The obtained metric values demonstrate that Go implements a deliberately simplified object model focused on composition, interface-based abstraction, and implementation minimalism rather than on classical object-oriented extensibility. The relatively high value of $TSSI(\text{Golang}) \approx 0.6334$ indicates that the language partially separates subtyping from inheritance through its interface system. Go does not support implementation inheritance, while interfaces are satisfied implicitly, which creates a form of structural subtyping independent from class hierarchies. Such separation corresponds to the distinction between inheritance and subtype polymorphism discussed by Cardelli and Wegner \cite{cardelli1985understanding}. At the same time, the absence of variance formalization and the intentionally restricted generic system reduce the overall expressiveness of the type model. The relatively high value of $ISRI(\text{Golang}) \approx 0.7167$ is primarily explained by the absence of multiple inheritance and by deterministic method resolution semantics. Since Go eliminates inheritance hierarchies entirely, it avoids a substantial portion of ambiguity traditionally associated with object-oriented inheritance systems and the fragile base class problem. However, this simplification is achieved by reducing inheritance expressiveness rather than by introducing advanced hierarchy management mechanisms.

The metrics related to composition, polymorphism, and formal correctness demonstrate the limitations of Go as a fully object-oriented language. The low values of $CRFI(\text{Golang}) = 0.25$ and $PDEI(\text{Golang}) = 0.245$ show that the language provides only minimal support for reusable behavioral composition and advanced dispatch mechanisms. Go lacks traits, mixins, multiple dispatch, and native delegation constructs, relying almost exclusively on interfaces and explicit composition patterns. This design reflects the Go philosophy of simplicity and explicitness introduced by Pike et al. in the original language design discussions \cite{pike2012go}. Similarly, the low values of $FSVI(\text{Golang}) = 0.275$ and $SCSI(\text{Golang}) = 0.3$ indicate that Go prioritizes pragmatic concurrency support over formal verification guarantees. Although the language includes high-level concurrency primitives, it does not provide static race-freedom guarantees, ownership-based aliasing control, or native contract mechanisms. The moderate value of $ORCI(\text{Golang}) \approx 0.5917$ additionally reflects that Go supports several object construction styles and runtime introspection through reflection, but intentionally restricts metaobject extensibility and runtime intercession. Overall, the collected metrics confirm that Go represents a minimalistic and composition-oriented reinterpretation of object-oriented programming rather than a feature-rich object-oriented model.

\subsection{Comparison summary: Scala}

The obtained results demonstrate that Scala provides one of the most formally expressive object models among contemporary mainstream object-oriented languages. The high value of $TSSI=0.8$ indicates that the language combines nominal and structural typing while also supporting formally specified variance annotations for parametric polymorphism. Such a result reflects Scala’s orientation toward mathematically rigorous type safety and flexible subtype relations. The language specification explicitly separates several subtype mechanisms from direct implementation inheritance through traits, abstract type members, and structural types, although this separation is not entirely independent semantically, which explains the partial reduction of the $S_{rel}$ component. This observation is consistent with the description of Scala’s type system provided by Odersky et al., where the language is characterized as a synthesis of object-oriented and functional abstraction mechanisms with a highly expressive static type system \cite{odersky2004overview}. In addition, the relatively high $ISRI=0.8167$ confirms that Scala introduces strong hierarchy-management mechanisms, including deterministic linearization, explicit override semantics, and restrictions aimed at reducing fragile base class effects. The trait linearization model follows deterministic composition semantics similar to those discussed in trait-oriented research by Schärli et al. \cite{scharli2003traits}.

The results also show that Scala strongly emphasizes compositional reuse and abstraction flexibility. The value $CRFI=0.6875$ reflects extensive support for trait-based composition and interface-oriented behavioral reuse, which substantially reduces dependence on classical inheritance hierarchies. Traits in Scala act not only as contracts but also as reusable behavioral units, corresponding to the compositional approach proposed in trait research \cite{scharli2003traits}. Similarly, the relatively high $ORCI=0.6917$ demonstrates support for several object creation paradigms and a fully reference-oriented object model. However, the moderate $PDEI=0.62$ indicates that Scala remains largely based on single dispatch despite advanced pattern matching facilities. More significant limitations become visible in the areas of encapsulation enforcement and formal verification. The low values of $AE=0.3381$ and $FSVI=0.275$ are primarily caused by the absence of strict runtime encapsulation guarantees and the lack of native Design by Contract constructs. While Scala provides expressive abstraction mechanisms, reflection and JVM interoperability weaken strict encapsulation boundaries. Furthermore, formal correctness specifications rely mostly on external libraries rather than language-level constructs. Finally, the moderate $SCSI=0.4$ demonstrates that Scala only partially addresses immutability and concurrency safety at the language level: although immutable collections and actor-based concurrency models are available, the language does not enforce race freedom or immutability by default. This partially confirms the broader observation that advanced abstraction capabilities in hybrid object-functional languages do not automatically imply strong guarantees of state safety or formal verifiability \cite{clarke1998ownership,meyer2002applying}.

\subsection{Comparison summary: JavaScript}

The obtained metric values demonstrate that JavaScript provides only limited support for classical object-oriented semantics despite offering object-based programming constructs. The low value of $TSSI=0.2$ reflects the absence of a formally specified static type system and the lack of explicit subtype relations in the language specification. As noted by Cardelli, subtype reasoning requires formally defined substitutability relations independent from implementation reuse \cite{cardelli1985understanding}, whereas JavaScript historically relies on dynamic prototype delegation instead of explicit subtyping. Similarly, the relatively weak values of $AE=0.3071$ and $FSVI=0.275$ indicate insufficient language-level support for strict encapsulation and formal specification. Although ECMAScript modules provide explicit export boundaries, the language still exposes reflective runtime mechanisms and lacks native Design by Contract constructs. Meyer emphasized that reliable object-oriented abstraction requires explicit contractual semantics and controlled visibility mechanisms \cite{meyer1997object,meyer2002applying}, which JavaScript only partially achieves through conventions and runtime checks. The extremely low values of $CRFI=0.075$, $PDEI=0.07$, and $SCSI=0.1$ further demonstrate that JavaScript prioritizes flexibility and dynamic adaptation over formal correctness guarantees. In particular, the language lacks native trait systems, multiple dispatch, static exhaustiveness guarantees, and compile-time race prevention mechanisms.

At the same time, JavaScript achieves one of the highest values for object representation flexibility, with $ORCI=0.925$. This result reflects the dynamic and highly extensible nature of the ECMAScript object model. The language supports multiple object creation paradigms, including constructor-based instantiation, object literals, and prototype cloning, while also exposing powerful meta-programming facilities through reflective APIs and prototype manipulation. Such characteristics correspond to the metaobject-oriented ideas described by Kiczales et al. \cite{kiczales1991art}. The moderate value of $ISRI=0.4917$ additionally indicates that JavaScript avoids some classical inheritance complexity problems due to its prototype-based model and deterministic property lookup semantics, but simultaneously lacks formal safeguards against fragile base hierarchy effects. Overall, the metric profile confirms conclusions commonly discussed in research on dynamic languages: JavaScript provides high runtime flexibility and extensibility, but achieves this at the cost of weak formal guarantees, limited encapsulation safety, and reduced suitability for constructing large-scale verifiable object-oriented systems.

\subsection{Comparison summary: Python}

The obtained metric values demonstrate that Python provides a comparatively flexible and expressive object model, but this flexibility is achieved at the expense of strict semantic guarantees and formal safety mechanisms. The high value of the Type System and Subtyping Soundness Index ($TSSI = 0.8583$) reflects Python’s hybrid typing approach, which combines nominal inheritance with structural typing through protocols introduced in PEP~544. This confirms the observation of Cardelli that structural and nominal subtyping solve different compatibility problems and may coexist within a single type system~\cite{cardelli1985understanding}. At the same time, the relatively high Object Representation and Creation Index ($ORCI = 0.85$) indicates the expressive meta-object capabilities of Python, including runtime reflection and dynamic object manipulation, which correspond to the principles of reflective systems described by Kiczales et al.~\cite{kiczales1991art}. The moderate values of the Abstraction and Encapsulation score ($AE = 0.5207$) and Composition and Reuse Flexibility Index ($CRFI = 0.5125$) show that Python supports interfaces, mixin-oriented reuse, and delegation patterns, but lacks strong encapsulation barriers and formalized trait composition semantics. This observation is consistent with research on traits and compositional reuse, where explicit conflict resolution and formal requirements are considered necessary for scalable composition mechanisms~\cite{scharli2003traits}.

However, the remaining metrics reveal substantial limitations of Python from the perspective of robust object-oriented system design. The low value of the Inheritance Semantics Robustness Index ($ISRI = 0.3417$) reflects the complexity introduced by unrestricted multiple inheritance and the absence of strong compile-time hierarchy constraints. Although Python formally specifies deterministic method resolution order through C3 linearization, it provides only limited protection against fragile base class problems and accidental overriding. The extremely low value of the Polymorphism and Dispatch Efficiency Index ($PDEI = 0.07$) demonstrates that Python relies almost entirely on conventional single dispatch without compile-time dispatch verification or optimization-oriented hierarchy restrictions. Similarly, the low State and Concurrency Safety Index ($SCSI = 0.1$) confirms the absence of language-level guarantees for immutability, aliasing control, or data-race freedom, forcing correctness to depend primarily on programmer discipline rather than static enforcement mechanisms. Finally, the Formal Specification and Verifiability Index ($FSVI = 0.275$) indicates that Python does not provide native Design by Contract facilities and relies mostly on external libraries or runtime assertions. This contrasts with the Design by Contract approach proposed by Meyer, where contracts are treated as first-class semantic elements of the language specification~\cite{meyer2002applying}. Overall, the results characterize Python as a language optimized for flexibility, runtime adaptability, and rapid software evolution rather than for formally verifiable and semantically constrained object-oriented architecture.

\subsection{Comparison summary: Ruby}

The obtained metric values demonstrate that Ruby prioritizes runtime flexibility and metaprogramming expressiveness over static correctness guarantees and formally constrained object-oriented semantics. The low value of the TSSI metric ($0.15$) indicates that Ruby lacks a formally expressive subtype system: the language follows a dynamically typed object model without explicit variance semantics or formal separation between inheritance and subtyping. This observation corresponds to the classical distinction between inheritance and subtype polymorphism discussed by Cardelli and Wegner, who emphasize that dynamic object systems often conflate implementation reuse with substitutability \cite{cardelli1985understanding}. Similarly, the low AE value ($0.2107$) reflects Ruby’s intentionally permissive encapsulation model. Although the language provides several visibility modifiers (\texttt{private}, \texttt{protected}, \texttt{public}), these mechanisms remain advisory rather than strictly enforced, since reflection and metaprogramming allow unrestricted runtime access to object internals. Such behavior is consistent with the reflective philosophy described in metaobject protocol research \cite{kiczales1991art}. The absence of strict modular isolation and native contract mechanisms further reduces the FSVI score ($0.275$), as Ruby relies primarily on runtime assertions and external libraries instead of language-level specification constructs.

At the same time, Ruby demonstrates comparatively strong results in flexibility-oriented metrics. The high ORCI value ($0.925$) reflects the language’s highly dynamic object representation model, extensive metaprogramming facilities, and support for multiple object creation paradigms. Ruby’s object system allows runtime structural modification, singleton methods, and dynamic class extension, closely aligning with reflective object model approaches described in \cite{kiczales1991art}. The CRFI metric ($0.5125$) also indicates moderate support for composition-oriented reuse through modules and mixins, which partially mitigate limitations of single inheritance. However, inheritance-related robustness remains limited despite deterministic method lookup semantics. The ISRI value ($0.4917$) shows that Ruby benefits from single inheritance simplicity and deterministic linearization, but lacks stronger hierarchy restriction mechanisms such as sealed hierarchies, mandatory override annotations, or non-virtual-by-default semantics. Particularly significant are the extremely low PDEI ($0.07$) and SCSI ($0$) values, which reveal the absence of static dispatch safety guarantees, exhaustive polymorphism checking, ownership constraints, immutability guarantees, and compile-time concurrency safety mechanisms. This confirms that Ruby’s object model is optimized primarily for developer productivity and runtime adaptability rather than for formal correctness, static verifiability, or safe large-scale composition.

\subsection{Comparison summary: Zonnon}

The obtained metric values for the Zonnon language demonstrate a heterogeneous level of maturity of the object model. The strongest side of the language is the abstraction and encapsulation subsystem ($AE\approx 0.8232$), which indicates well-developed support for modularity, clear separation of interface and implementation, as well as strict access control. This result is consistent with the classical principles of modular design~\cite{parnas1972criteria}, as well as with the concept of encapsulation as a key mechanism for managing complexity in OOP~\cite{meyer1997object}. Additionally, the high ISRI score ($\approx 0.7667$) indicates a well-formalized inheritance semantics, including a deterministic method resolution order and partial hierarchy protection mechanisms. This confirms the thesis about the need for a strict distinction between reuse and subtyping mechanisms, formulated in the work~\cite{cook1989inheritance}, which notes that inheritance and subtyping have different semantics, and mixing these concepts leads to architectural limitations.

At the same time, a number of key aspects of the object model are implemented significantly less well. Low values of CRFI ($0.25$) and SCSI ($\approx 0.3667$) indicate limited support for composition and insufficient guarantees of state security and competitiveness. The lack of expressive traits/mixins mechanisms and automatic delegation reduces the flexibility of reuse, which contradicts modern approaches that consider composition as a preferred alternative to inheritance~\cite{scharli2003traits}. Similarly, weak guarantees of immutability and link control confirm the findings of ownership and aliasing studies, which emphasize that the absence of strict restrictions leads to competitive access errors~\cite{clarke1998ownership}. The moderate values of TSSI ($0.55$), ORCI ($\approx 0.4417$), and PDEI ($0.445$) additionally show that the language only partially implements the formal aspects of the model system, polymorphism, and object model. In particular, the lack of formalization of variation and limited dispatch mechanisms contradict the results of fundamental research of standard systems, where strict formalization is a necessary condition for security~\cite{pierce2002types}. Taken together, the results show that Zonnon demonstrates a strong architectural discipline at the level of modularity and inheritance, but is inferior to modern languages in composition flexibility, formal rigor of the standard system and guarantees of execution safety, which confirms the need to develop a more balanced object model within the framework of the proposed language.
